% CLAUDE REPORT::DERIVBRIDGE-0001::2026-11-11T09:55-08:00
% Compile with XeLaTeX or pdfLaTeX. P versus NP remains OPEN.
\documentclass[11pt]{article}
\usepackage[margin=1in]{geometry}
\usepackage{amsmath,amssymb,amsthm}
\usepackage{hyperref}
\newtheorem{theorem}{Theorem}
\newtheorem{proposition}[theorem]{Proposition}
\newtheorem{corollary}[theorem]{Corollary}
\newcommand{\Wb}{\overline W}
\title{CLAUDE DERIVBRIDGE-0001: where the bridge $\mathrm{SAT}\in\mathrm P\Rightarrow\Wb_R\le\mathrm{poly}$ sits\\
\large world-ledger rows L41--L49 for the worst-case minimum-derivation program}
\author{Claude routine (scheduled run), repository \texttt{j15marti/pnp}}
\date{November 11, 2026, 09:55 PST}
\begin{document}
\maketitle

\noindent\textbf{STATUS.} Propositions 1, 4 and Theorem 2: PROVED/EXACT (elementary).
Corollary 3: CONDITIONAL on the corpus mutual $p$-simulation of $R_{\rm mix341}$ and
Extended Resolution. Pich--Santhanam row L48: KNOWN THEOREM (abstract read; full
statement audited in the corpus by LSPV-1/406). Ledger build: COMPUTATIONAL
(49/49 rows classified, 0 consistency conflicts). B2, B3 for any system without a
proved lower bound, and P versus NP: \textbf{OPEN}.
This run adds no new lower bound and no new bridge. It sorts the shared target's
statements by what they are worth.

\section{Shared target and the rationale for free optimal derivations}
For a fixed sound, complete, polynomial-time-checkable refutation system $R$ with nonnegative
node weights,
\[
\Wb_R(N)=\max_{F\in\mathrm{UNSAT},\ |\langle F\rangle|\le N}\ \min_{\pi:F\vdash_R\bot}\operatorname{cost}_R(\pi).
\]
The inner minimum gives away perfect knowledge of the best reduction sequence. That is
legitimate for a lower bound: any executor of a particular derivation pays at least the
minimum, and search can only add work. We assume throughout that
$\operatorname{cost}_R$ and the bit length $|\pi|$ are polynomially related,
$|\pi|\le q(\operatorname{cost}_R(\pi)+N)$ and $\operatorname{cost}_R(\pi)\le q(|\pi|)$,
as in the corpus' bounded-arity node grammars (REFEREE334 shows why this must be stated).
Write $L_R(F)$ for the minimum proof length; then $\Wb_R$ is polynomial iff $R$ is
polynomially bounded ($p$-bounded).

\section{Results}
\begin{proposition}[Proved lower bound makes B3 a restatement]\label{p1}
If $\Wb_R$ is unconditionally superpolynomial, then
$\mathrm{B3}(R)\iff\mathrm{SAT}\notin\mathrm P$. Consequently ``there is a fixed $R$ with
$\Wb_R$ superpolynomial and $\mathrm{B3}(R)$'' is equivalent to $\mathrm P\ne\mathrm{NP}$.
\end{proposition}
\begin{proof}
An implication with a false consequent is equivalent to the negation of its hypothesis.
For the second claim: ($\Leftarrow$) take $R=$ Resolution (Haken 1985), whose B3 holds
vacuously when $\mathrm{SAT}\notin\mathrm P$; ($\Rightarrow$) B3 plus $\mathrm{SAT}\in\mathrm P$
would make $\Wb_R$ polynomial.
\end{proof}
So for Resolution, bounded-depth Frege, Polynomial Calculus or Cutting Planes on their
known hard families, ``proving the bridge'' is literally proving $\mathrm P\ne\mathrm{NP}$
(same point as LSPV-1/398). The two-step program makes progress only if B3 is proved
for a system \emph{before} its lower bound is known.

\begin{theorem}[Optimality form of B3]\label{t2}
For every Cook--Reckhow UNSAT system $R$ as above,
\[
\mathrm{B3}(R)\iff\bigl[\mathrm P=\mathrm{NP}\Rightarrow R\text{ is size-optimal}\bigr],
\]
where size-optimal means: for every Cook--Reckhow system $Q$ there is a polynomial $r$
with $L_R(F)\le r(L_Q(F)+|F|)$ for all UNSAT $F$.
\end{theorem}
\begin{proof}
Assume $\mathrm P=\mathrm{NP}$. If $R$ is $p$-bounded then $L_R(F)\le r(|F|)\le r(L_Q(F)+|F|)$
for every $Q$. Conversely, let $A$ be a polynomial-time UNSAT decider and let $Q_A$ accept
$(F,\varepsilon)$ iff $A$ rejects $F$; $Q_A$ is a Cook--Reckhow system with $L_{Q_A}(F)=0$.
Optimality gives $L_R(F)\le r(|F|)$, so $R$ is $p$-bounded. Under $\mathrm P\ne\mathrm{NP}$ both
sides hold vacuously.
\end{proof}
The theorem relativizes as an implication. It does not rescue a relativizing proof of B3
itself (REFEREE296's barrier stands).

\begin{corollary}[B3 moves up the simulation order]\label{c3}
If $R\le R'$ in size ($L_{R'}\le\mathrm{poly}(L_R+|F|)$), then
$\mathrm{B3}(R)\Rightarrow\mathrm{B3}(R')$. In particular
$\mathrm{B3}(\text{Frege})\Rightarrow\mathrm{B3}(\mathrm{EF})$, and, CONDITIONAL on the corpus
mutual simulation $R_{\rm mix341}\equiv_p\mathrm{ER}\equiv_p\mathrm{EF}$ (REFEREE350's explicit root-only
completion; LSPV-1/406 notes that the native LOCAL-CONVERT grammar is still underspecified),
$\mathrm{B3}(R_{\rm mix341})\iff\mathrm{B3}(\mathrm{EF})$.
\end{corollary}
\begin{proof} $p$-boundedness transfers upward along size simulation. \end{proof}

So the project's B3 for its fixed calculus is the long-open problem of whether
$\mathrm P=\mathrm{NP}$ forces Extended Frege to be $p$-bounded. Equivalently, by Theorem
\ref{t2}: does a superpolynomial EF lower bound imply $\mathrm P\ne\mathrm{NP}$?
The two obligations pull in opposite directions. B3 gets easier as $R$ gets stronger
(Corollary \ref{c3}), and lower bounds get harder. For the strongest natural choice, $R_T$
(a proof is a $T$-proof of the sentence $\mathrm{UNSAT}(\ulcorner F\urcorner)$ for a consistent
r.e.\ $T\supseteq\mathrm{PV}$, which is sound by $\Sigma_1$-completeness), B3 follows from
``$\mathrm P=\mathrm{NP}$ via some decider whose rejection soundness $T$ proves''. This is the
corpus' one-sided PV route, and no lower bound for $R_T$ is in sight.

\begin{proposition}[Worst-case B2 for Resolution is ETH-type]\label{p4}
$\mathrm{B2}(\mathrm{Res})$ implies that every correct SAT decider $A$ has
$T^{\rm wc}_A(N)\ge 2^{N^{1/3-o(1)}}$ for infinitely many $N$.
\end{proposition}
\begin{proof}
Canonical $\mathrm{PHP}^{h+1}_h$ has encoded size $N=\Theta(h^3\log h)$ and needs $2^{\Omega(h)}$
resolution steps (Haken). So $\Wb_{\rm Res}(N)\ge 2^{N^{1/3-o(1)}}$, and
$\Wb\le p_A(N+T^{\rm wc}_A)$ forces $T^{\rm wc}_A\ge \Wb^{\Omega(1)}$.
\end{proof}
Hence B2 for any system with a proved exponential lower bound is strictly beyond
$\mathrm P\ne\mathrm{NP}$ (ledger L47, ``Implies $\mathrm P\ne\mathrm{NP}$ (open)'').

\section{Ledger rows added (L41--L49)}
\begin{center}\small
\begin{tabular}{lll}
Row & Statement & Category (auto)\\\hline
L41 & $\Wb_{\rm Res}$ superpolynomial (Haken; Urquhart) & Theorem\\
L42 & B3(Res) & Equivalent to $\mathrm P\ne\mathrm{NP}$\\
L43 & B3(EF) $\equiv$ B3($R_{\rm mix341}$) (conditional) & Half-win (true if $\mathrm P\ne\mathrm{NP}$)\\
L44 & EF $p$-bounded & Undetermined by P vs NP\\
L45 & $\exists R$: $\Wb_R$ superpolynomial $\wedge$ B3$(R)$ & Equivalent to $\mathrm P\ne\mathrm{NP}$\\
L46 & B1(Res), pointwise & Refuted in every world\\
L47 & B2(Res), worst case & Implies $\mathrm P\ne\mathrm{NP}$ (open)\\
L48 & Pich--Santhanam feasible-antichecker theorem & Theorem\\
L49 & Theorem \ref{t2} & Theorem\\
\end{tabular}
\end{center}
Regenerate with \texttt{python research/agents/claude/world-ledger/build\_ledger.py}
(requires openpyxl, runs in about 1\,s). Output: 49 rows, all consistency checks OK.

\section{Literature check (this run)}
\begin{itemize}
\item J.~Pich, R.~Santhanam, \emph{Towards $\mathrm P\ne\mathrm{NP}$ from Extended Frege lower bounds},
arXiv:2312.08163 / ECCC TR23-199 (abstract read). Under $S^1_2$-provability hypotheses
(feasible anticheckers; witnessing $\mathrm{NP}\not\subseteq\mathrm{P/poly}$; or condition I plus an
OWF-to-SAT reduction), \emph{EF not $p$-bounded $\Rightarrow\mathrm P\ne\mathrm{NP}$}. This is a
conditional proof of the contrapositive of B3(EF), which is the clue most directly relevant to the
current gap. The open part is the provability hypothesis. LSPV-1/406 already carried
Theorem 10 into $R_{\rm mix341}$; this run adds its place in the ledger.
\item Arteche, Khaniki, Pich, Santhanam, ICALP 2024 (arXiv:2405.02232), from search results:
an iEF version giving $\#\mathrm P\not\subseteq\mathrm{FP/poly}$ from iEF lower bounds, under
iEF-provable derandomization. The same group observes that establishing the
proof-complexity-to-circuit connection unconditionally would already give
$\mathrm{NEXP}\not\subseteq\mathrm{P/poly}$. That is a warning: an unconditional B3(EF)-type
transfer would carry consequences at least of that strength. CANDIDATE barrier; not verified at
the theorem level this run.
\item Cook--Reckhow 1979 ($\exists$ $p$-bounded system $\iff\mathrm{NP}=\mathrm{coNP}$);
Kraj\'{\i}\v{c}ek--Pudl\'ak 1989 ($\mathrm{NE}=\mathrm{coNE}\Rightarrow$ optimal systems exist;
$\mathrm E=\mathrm{NE}\Rightarrow p$-optimal): KNOWN THEOREMS, used only through Theorem \ref{t2}.
\end{itemize}
Transfer classification. Bounded arithmetic to propositional proofs (Cook's PV to EF) is a
DIRECTLY TRANSFERABLE THEOREM, but only under a provability hypothesis. ``P = NP gives a
short proof system'' is a FAILED TRANSFER to a \emph{fixed} $R$ unless optimality is proved
(Theorem \ref{t2}). Post-quantum ledger: no item applies to this purely logical question
(FAILED TRANSFER of items 1--4; item 5, black-box extraction, falls under REFEREE296's
relativization barrier).

\section*{Corpus reviewed}
AGENTS.md (blob 4ae90288); root \texttt{main.tex} (blob ee29725d, REFEREE295/296 executive
blocks); the 2026-11-02 five-role worst-case refocus directive, with LSPV-2/341, LSPV-4/341,
LSPV-3/342 and REFEREE295 appended; LSPV-1/336 (clocked-optimality reflection), LSPV-1/398
(Resolution B3 firewall), LSPV-1/406 (feasible anticheckers); repository-wide greps for PV,
reflection, $p$-simulation, optimal proof system and antichecker (17 files); the Claude
world-ledger and earlier CLAUDE reports. Head at start: \texttt{bd1389dd}.

\section*{Duplication / contradiction audit}
Proposition \ref{p1} duplicates LSPV-1/398's corrected label, credited. Theorem \ref{t2} is
close to LSPV-1/336 and REFEREE296 (optimality $\Rightarrow$ B1); the converse direction
under $\mathrm P=\mathrm{NP}$ and the ``iff'' form were not found stated in the corpus. Corollary
\ref{c3}'s identification with EF follows from REFEREE350 plus the upward-monotonicity remark.
The ledger rows are new. No contradictions found.

\section*{What changed relative to root main}
Nothing in \texttt{main.tex}. This routine owns no ownership region, so the report and ledger
are offered to LSPV-5 as evidence. B3 for the fixed calculus is now located as an instance of
a named open problem, and the corresponding ledger categories are recorded.

\section*{Adopted / rejected / retired prior components}
Adopted: the definition of $\Wb_R$; B1/B2/B3; REFEREE296's relativization barrier;
Pich--Santhanam Theorem 10 as audited by LSPV-1/406. Rejected: reading Cook--Reckhow as B3
for a fixed $R$ (already rejected in the corpus). Retired: none.

\section*{Stale or overstated claims found}
None new. Reminder: Haken-type lower bounds (L41) are true in every world. They become
evidence about machines only through B3, and for those systems B3 is P $\ne$ NP itself (L42).

\section*{Next missing implication}
Either (a) prove the $S^1_2$ feasible-antichecker hypothesis of L48, which would give
B3(EF), hence B3($R_{\rm mix341}$) conditional on Corollary \ref{c3}; or (b) find a
system strictly weaker than EF with an EF-free proof that $\mathrm P=\mathrm{NP}$ makes it
optimal. Both are OPEN. P versus NP remains OPEN.

% END_OF_VERIFIED_REPORT::CLAUDE::DERIVBRIDGE-0001
\end{document}
